\documentclass[10pt,a4paper]{amsart}
\usepackage{amsmath,amssymb,amsthm,graphicx}
\usepackage{verbatim}
\newenvironment{reference}{\comment}{\endcomment}
%\newenvironment{reference}{}{}
\newenvironment{slogan}{\comment}{\endcomment}
\newenvironment{history}{\comment}{\endcomment}

% For commutative diagrams we use Xy-pic
\usepackage[all]{xy}

% We use 2cell for 2-commutative diagrams.
\xyoption{2cell}
% Theorem environments.
%
\theoremstyle{plain}
\newtheorem{theorem}[subsection]{Theorem}
\newtheorem{proposition}[subsection]{Proposition}
\newtheorem{lemma}[subsection]{Lemma}

\theoremstyle{definition}
\newtheorem{definition}[subsection]{Definition}
\newtheorem{example}[subsection]{Example}
\newtheorem{exercise}[subsection]{Exercise}
\newtheorem{situation}[subsection]{Situation}

\theoremstyle{remark}
\newtheorem{remark}[subsection]{Remark}
\newtheorem{remarks}[subsection]{Remarks}

\numberwithin{equation}{subsection}

% Macros
%
\def\lim{\mathop{\mathrm{lim}}\nolimits}
\def\colim{\mathop{\mathrm{colim}}\nolimits}
\def\Spec{\mathop{\mathrm{Spec}}}
\def\Hom{\mathop{\mathrm{Hom}}\nolimits}
\def\Ext{\mathop{\mathrm{Ext}}\nolimits}
\def\SheafHom{\mathop{\mathcal{H}\!\mathit{om}}\nolimits}
\def\SheafExt{\mathop{\mathcal{E}\!\mathit{xt}}\nolimits}
\def\Sch{\mathit{Sch}}
\def\Mor{\mathop{\mathrm{Mor}}\nolimits}
\def\Ob{\mathop{\mathrm{Ob}}\nolimits}
\def\Sh{\mathop{\mathit{Sh}}\nolimits}
\def\NL{\mathop{N\!L}\nolimits}
\def\CH{\mathop{\mathrm{CH}}\nolimits}
\def\proetale{{pro\text{-}\acute{e}tale}}
\def\etale{{\acute{e}tale}}
\def\QCoh{\mathit{QCoh}}
\def\Ker{\mathop{\mathrm{Ker}}}
\def\Im{\mathop{\mathrm{Im}}}
\def\Coker{\mathop{\mathrm{Coker}}}
\def\Coim{\mathop{\mathrm{Coim}}}

% Boxtimes
%
\DeclareMathSymbol{\boxtimes}{\mathbin}{AMSa}{"02}

%
% Macros for moduli stacks/spaces
%
\def\QCohstack{\mathcal{QC}\!\mathit{oh}}
\def\Cohstack{\mathcal{C}\!\mathit{oh}}
\def\Spacesstack{\mathcal{S}\!\mathit{paces}}
\def\Quotfunctor{\mathrm{Quot}}
\def\Hilbfunctor{\mathrm{Hilb}}
\def\Curvesstack{\mathcal{C}\!\mathit{urves}}
\def\Polarizedstack{\mathcal{P}\!\mathit{olarized}}
\def\Complexesstack{\mathcal{C}\!\mathit{omplexes}}
% \Pic is the operator that assigns to X its picard group, usage \Pic(X)
% \Picardstack_{X/B} denotes the Picard stack of X over B
% \Picardfunctor_{X/B} denotes the Picard functor of X over B
\def\Pic{\mathop{\mathrm{Pic}}\nolimits}
\def\Picardstack{\mathcal{P}\!\mathit{ic}}
\def\Picardfunctor{\mathrm{Pic}}
\def\Deformationcategory{\mathcal{D}\!\mathit{ef}}
\usepackage{cleveref}
\crefname{originallemma}{original lemma}{original lemmas}
\Crefname{originallemma}{Original lemma}{Original lemmas}
\makeatletter
\newlabel{categories-lemma-tensor-dual}{{0FFS}{source}{}{sourcepointer-1}{}}
\newlabel{categories-lemma-tensor-dual@cref}{{[originallemma][0][]0FFS}{[1][0][]source}{}{}{}}
\makeatother
\crefname{originallemma}{original lemma}{original lemmas}
\Crefname{originallemma}{Original lemma}{Original lemmas}
\makeatletter
\newlabel{homology-lemma-Karoubian-dual}{{0FFU}{source}{}{sourcepointer-2}{}}
\newlabel{homology-lemma-Karoubian-dual@cref}{{[originallemma][0][]0FFU}{[1][0][]source}{}{}{}}
\makeatother
\crefname{originallemma}{original lemma}{original lemmas}
\Crefname{originallemma}{Original lemma}{Original lemmas}
\makeatletter
\newlabel{homology-lemma-additive-dual}{{0FFT}{source}{}{sourcepointer-3}{}}
\newlabel{homology-lemma-additive-dual@cref}{{[originallemma][0][]0FFT}{[1][0][]source}{}{}{}}
\makeatother
\crefname{originalreference}{original reference}{original references}
\Crefname{originalreference}{Original reference}{Original references}
\makeatletter
\newlabel{lemma-dual}{{08I1}{source}{}{sourcepointer-4}{}}
\newlabel{lemma-dual@cref}{{[originalreference][0][]08I1}{[1][0][]source}{}{}{}}
\makeatother
\title{Every motive has a left dual}
\author{The Stacks Project Authors}
\date{stacks-project:weil.tex:local-snapshot-20261012; GNU FDL 1.2 or later}
\begin{document}
\maketitle
\noindent\textit{Editorial note.} All delivered mathematical text is the author's own, copied from the pinned Stacks Project chapter file weil.tex; no statement or proof step was rewritten, reordered or removed.\par\medskip
\noindent\textit{Editorial note.} The counted proof is the author's complete argument: the reduction to twisted motives through summands, the removal of the twist, and the reduction to irreducible components.\par\medskip
\noindent\textit{Editorial note.} The label lemma-dual-general names the counted statement, and the proof environment immediately following it in the source is delivered verbatim as its proof.\par\medskip
\noindent\textit{Editorial note.} The counted statement is carried by the author's own lemma environment in the chapter source.\par\medskip
\noindent\textit{Editorial note.} Pre-open gate: python operations/agent-authoring/tools/precard\_gate.py --source weil.tex --label lemma-dual-general --start 1237 --end 1256 --local-id 28160 returned RESULT=OK (L28160 free, occupancy FREE, public interval clear). Pre-assembly tag self-check: canonical tag resolved and not present in reference/public-items.jsonl.\par\medskip
\noindent\textit{Editorial note.} The statement and proof are copied from the pinned chapter file; the cited inputs are kept as external pointers with their canonical Stacks tags.\par\medskip
\noindent\textit{Allowed standard prerequisites (editorial).}\begin{itemize}
\item Summands of objects with a dual have a dual. Homology, the lemma that a summand of an object admitting a left dual again admits one. Reduces the statement from an arbitrary motive to a twisted motive of a smooth projective scheme.
\item Twisting an object providing a dual. The construction of the twist as the tensor product with the Tate object, together with the categorical lemma that tensoring with an object admitting a dual preserves the existence of duals. Reduces the statement from a twisted motive to the untwisted motive of the scheme.
\item Additivity of duals and duality of the motive of an equidimensional scheme. Homology, the lemma on duals of direct sums, and Same chapter, the lemma exhibiting the twisted motive of an equidimensional smooth projective scheme as a left dual of its motive. Reduces to the irreducible components, where the duality statement applies.
\end{itemize}\medskip
\ifdefined\hypertarget\hypertarget{sourcepointer-1}{}\fi
\noindent\textit{Source reference.} Original lemma 0FFS: sources/stacks-chow/tags.txt, source lines 19297--19297. Prerequisite pointer invoked at proof line(s) 1250; it names a standard fact of a sibling chapter that the counted argument uses, and every step of the counted proof built on it is retained verbatim in the delivered body.\par\medskip
\ifdefined\hypertarget\hypertarget{sourcepointer-2}{}\fi
\noindent\textit{Source reference.} Original lemma 0FFU: sources/stacks-chow/tags.txt, source lines 19299--19299. Prerequisite pointer invoked at proof line(s) 1245; it names a standard fact of a sibling chapter that the counted argument uses, and every step of the counted proof built on it is retained verbatim in the delivered body.\par\medskip
\ifdefined\hypertarget\hypertarget{sourcepointer-3}{}\fi
\noindent\textit{Source reference.} Original lemma 0FFT: sources/stacks-chow/tags.txt, source lines 19298--19298. Prerequisite pointer invoked at proof line(s) 1254; it names a standard fact of a sibling chapter that the counted argument uses, and every step of the counted proof built on it is retained verbatim in the delivered body.\par\medskip
\ifdefined\hypertarget\hypertarget{sourcepointer-4}{}\fi
\noindent\textit{Source reference.} Original reference 08I1: sources/stacks-weil/pinned-original-source.tex, source lines 1186--1186. Prerequisite pointer invoked at proof line(s) 1255; it names a standard fact of the same chapter that the counted argument uses, and every step of the counted proof built on it is retained verbatim in the delivered body.\par\medskip
\begin{lemma}
\label{lemma-dual-general}
Every object of $M_k$ has a left dual.
\end{lemma}

\begin{proof}
Let $M = (X, p, m)$ be an object of $M_k$. Then $M$ is a summand of
$(X, 1, m) = h(X)(m)$.
By Homology, Lemma \ref{homology-lemma-Karoubian-dual}
it suffices to show that
$h(X)(m) = h(X) \otimes \mathbf{1}(m)$ has a dual.
By construction $\mathbf{1}(-m)$ is a left dual of $\mathbf{1}(m)$.
Hence it suffices to show that $h(X)$ has a left dual, see
Categories, Lemma \ref{categories-lemma-tensor-dual}.
Let $X = \coprod X_i$ be the decomposition of $X$ into
irreducible components. Then $h(X) = \bigoplus h(X_i)$
and it suffices to show that $h(X_i)$ has a left dual, see
Homology, Lemma \ref{homology-lemma-additive-dual}.
This follows from Lemma \ref{lemma-dual}.
\end{proof}

\end{document}
